\begin{textsample}{POS Dim 2 – es – Score 22.00 – es/es\_em\_13.txt}  \label{ex:f2_pos_243}
Estrutura de \textbf{Oferta} Os \textbf{Itinerários} Formativos do Estado do Espírito Santo foram construídos articulando as condições de \textbf{oferta} da rede estadual e o perfil dos estudantes.
Com o objetivo de potencializar o envolvimento e a participação ativa da comunidade escolar, foram realizadas diferentes ações, entre elas : a aplicação de questionários, a realização de rodas de conversas nas escolas e um encontro presencial reunindo representantes de estudantes de todo o estado[^3 ].
Essas ações tiveram o objetivo de compreender as expectativas e as necessidades dos estudantes da rede estadual e auxiliar as escolas e a Secretaria de Estado da Educação na promoção de melhorias que realmente tornem a escola mais significativa.
Se não bastasse, as pesquisas \textbf{diagnósticas} serviram para traçar \textbf{diretrizes} para a \textbf{oferta} dos \textbf{Itinerários} Formativos em toda a rede.
Esses \textbf{Itinerários} incluem a área de aprofundamento que o estudante poderá escolher e as unidades curriculares que serão comuns a todos e de participação obrigatória.
Dessa forma, nos próximos anos, as escolas da Rede Estadual deverão \textbf{ofertar} a \textbf{carga} \textbf{horária} mínima anual de 1.000h.
Ao final do Ensino Médio, os estudantes deverão concluir no mínimo 3.000h, distribuídas entre Formação Geral Básica ( 1.800h ) e \textbf{Itinerários} Formativos ( mínimo de 1.200h ), como descrito abaixo : - 1ª \textbf{Série} do Ensino Médio : 800h de Formação Geral Básica ( contemplando as quatro áreas do conhecimento ) e mínimo de 200h de \textbf{Itinerários} Formativos ( compostos por componentes integradores ) ; - 2ª \textbf{Série} do Ensino Médio : 600h de Formação Geral Básica ( contemplando as quatro áreas do conhecimento ) e mínimo de 400h de \textbf{Itinerários} Formativos ( compostos por componentes integradores e por aprofundamentos ) ; - 3ª \textbf{Série} do Ensino Médio : 400h de Formação Geral Básica ( contemplando as quatro áreas do conhecimento ) e 600h de \textbf{Itinerários} Formativos ( compostos por componentes integradores e por aprofundamentos ).
A \textbf{carga} \textbf{horária} do Ensino Médio, na parte \textbf{flexível}, será composta por Unidades Curriculares, que poderão ser obrigatórias ou \textbf{eletivas}, podendo o estudante optar por algumas das unidades \textbf{ofertadas} pelas escolas.
As Unidades Curriculares obrigatórias serão : Projeto de Vida, Estudo Orientado e aquelas do Aprofundamento que o estudante escolher.
Já as Unidades Curriculares de escolha dos estudantes serão as denominadas \textbf{Eletivas} e complementam os \textbf{Itinerários} Formativos.
As Unidades Curriculares foram contempladas por meio de oficinas, projetos, núcleos de estudo, módulos, incubadoras, clubes, entre outras possibilidades, de acordo com a proposta elaborada para cada um dos componentes propostos.
Essas definições estão presentes no \textbf{Guia} de Implementação do Novo Ensino Médio ( BRASIL, 2018e ), documento que apresenta as Unidades Curriculares como os elementos com \textbf{cargas} \textbf{horárias} pré-definidas e cujo objetivo é desenvolver competências específicas, seja da Formação Geral Básica, seja dos \textbf{Itinerários} Formativos.
O conjunto de Unidades Curriculares de um Itinerário deve desenvolver as habilidades de, pelo menos, um dos eixos estruturantes.
As aulas de Projeto de Vida poderão ser desenvolvidas por professores(as) de qualquer componente curricular, a partir de \textbf{diretrizes} encaminhadas pela Secretaria de Estado de Educação.
Elas têm como objetivo desenvolver competências socioemocionais que, apoiadas nos elementos cognitivos e nas experiências pessoais, promovam a consolidação de valores e conhecimentos na construção do projeto de vida do estudante.
As aulas de Estudo Orientado têm o objetivo de “ ensinar o estudante a estudar ”, proporcionando -lhe apoio e orientação em seus estudos diários, por meio de técnicas que o auxiliarão em seu processo de aprendizagem.
As aulas do componente “ \textbf{Eletivas} ” objetivam diversificar, aprofundar e/ou enriquecer os conteúdos e temas trabalhados nas disciplinas da Base Nacional Comum Curricular, considerando a interdisciplinaridade enquanto eixo metodológico.
As escolas deverão \textbf{ofertar} esta unidade curricular do Itinerário Formativo a partir de um \textbf{catálogo} de \textbf{eletivas} disponibilizado por esta Secretaria, podendo a escola também construí -las.
As \textbf{Eletivas} serão \textbf{ofertadas} pela escola, para a escolha dos estudantes, de acordo com o interesse e com o projeto de vida de cada um deles.
Junto com o Projeto de Vida, \textbf{Eletivas} e Estudo Orientado estão os percursos formativos de aprofundamento por área de conhecimento.
Esses aprofundamentos ocorrem nas duas últimas \textbf{séries} do Ensino Médio e são de escolha do estudante, dentro das possibilidades de \textbf{oferta} que a rede estadual irá oferecer.
A \textbf{carga} \textbf{horária} dos \textbf{Itinerários} Formativos é de no mínimo 1.200h e, desse período, pelo menos 800h são \textbf{destinadas} aos Aprofundamentos.
Alguns Aprofundamentos foram organizados, nesse primeiro momento de implementação do Novo \textbf{Currículo} do Espírito Santo, nas áreas de conhecimento e entre áreas, sendo um percurso em cada uma das duas últimas \textbf{séries} do Ensino Médio.
São inúmeras as possibilidades de arranjos para os Percursos de Aprofundamento, partindo sempre dos Referenciais Curriculares de 2018, elaborado pelo MEC.
Os \textbf{Itinerários} Formativos fomentam alternativas de diversificação e flexibilização, pelas unidades curriculares, de formatos ou formas de estudo e de atividades, estimulando a construção de percursos que \textbf{atendam} às características, interesses e necessidades dos estudantes e às demandas do meio social, privilegiando propostas com opções visando aos estudantes.
As propostas que seguem neste documento foram as primeiras planejadas pela \textbf{Equipe} de Implementação da Base Nacional Comum Curricular Etapa Ensino Médio, do Espírito Santo.
No entanto, se pretende que, posteriormente, as escolas tenham autonomia para criar os seus próprios Aprofundamentos, a partir dos Referenciais Curriculares legalizados pela \textbf{Portaria} nº 1.432 ( BRASIL, 2018c ).
De acordo com a Resolução nº 3/2018, § 10º, temos : > Formas diversificadas de \textbf{Itinerários} Formativos podem ser organizadas, desde que articuladas as dimensões do trabalho, da ciência, da tecnologia e da cultura, e definidas pela proposta pedagógica, \textbf{atendendo} necessidades, anseios e aspirações dos estudantes e a realidade da escola e seu meio. ( BRASIL, 2018b ).
Os percursos formativos de aprofundamentos podem ser organizados por diferentes arranjos, reunidos em uma ou mais áreas de conhecimento e, todos, articuladas nos eixos estruturantes.
Esses Aprofundamentos permitem que os estudantes possam fazer escolhas de acordo com seus interesses de formação.
Fazem referência a caminhos e estradas que devem ser organizados de acordo com as demandas regionais, em atendimento às demandas socioeconômicas e ambientais dos sujeitos e do mundo do trabalho.
A organização dos Aprofundamentos permite que uma escola centralize suas ações para determinadas áreas, de acordo com os eixos tecnológicos, otimizando recursos e aproveitando tecnologias comuns ( laboratórios e materiais ), bem como o quadro de professores e \textbf{técnicos} administrativos.
A apresentação dos Aprofundamentos conta com um planejamento geral, no qual estão estruturados os três módulos, que compõem o percurso de cada ano, e as diversas unidades curriculares que compõem cada módulo.
Para cada unidade curricular será apresentado um \textbf{detalhamento} que conta com as principais informações para o seu desenvolvimento.
As unidades curriculares têm duração anual e obedecerão aos critérios de avaliação trimestral.
O \textbf{detalhamento} das unidades curriculares traz a definição da(s) área(s) de conhecimento, dos eixos estruturantes, das habilidades relacionadas aos \textbf{Itinerários} Formativos associadas aos eixos, do tema, dos objetos de conhecimento, do tipo de unidades curriculares, da \textbf{carga} \textbf{horária}, do perfil docente, das possibilidades metodológicas e da avaliação.
Ao fim de cada Aprofundamento, pretende -se que os estudantes sejam capazes de alcançar sua formação humana para além de sua formação acadêmica, assim, é importante que cada Aprofundamento cuide da formação integral dos estudantes.
Dessa forma, para facilitar o entendimento destes itens, seguem abaixo as ementas do Estudo Orientado, da Disciplina \textbf{Eletiva} e do Projeto de Vida. [ ^3 ] : O questionário foi disponibilizado no mês de maio de 2019 e respondido na escola de forma individual, por todos os estudantes de 9º ano do Ensino Fundamental e 1ª \textbf{séries} do Ensino Médio.
Já as rodas de conversas foram realizadas nas escolas estaduais, com material orientador, e compilados pelas \textbf{Superintendências} Regionais de Educação.
Além disso, foi produzido pela Secretaria de Estado de Educação um encontro presencial, denominado “ IV Diálogos ”, com a participação de cerca de 900 representantes de estudantes, de todas as \textbf{superintendências} que compõem o Estado.

% matched lemmas: atender, carga, catálogo, currículo, destinar, detalhamento, diagnóstico, diretriz, eletivo, equipe, flexível, guia, horário, itinerário, oferta, ofertar, portaria, superintendência, série, técnico
\end{textsample}
